%%%%%%%%%%%%%%%%%%%%%%%%
% Package a rajouter   %
%%%%%%%%%%%%%%%%%%%%%%%%
\usepackage[centertags,intlimits]{amsmath}
\usepackage{amssymb,amsfonts}
\usepackage{mathrsfs}
\usepackage{bbm}
\usepackage{color}
\usepackage{xcolor}
\usepackage{tikz}
% Pie chart drawing library 
\usepackage{pgf-pie}  

%%%%%%%%%%%%%%%%%%%%%%%%%%%%%
% Macros specifiques (meca) %
%%%%%%%%%%%%%%%%%%%%%%%%%%%%%

%
% ENVIRONNEMENTS
%
\newcommand {\bea}{\begin{eqnarray*}}
\newcommand {\eea}{\end{eqnarray*}}
\newcommand {\dps}{\displaystyle}
%
% OPERATEURS TENSORIELS
%
\newcommand {\norm}[1]{\left|\!\left| {#1} \right|\!\right|} % norme de ...
\newcommand {\abs}[1]{\left|\:{#1}\:\right|} % valeur absolue de ...
\newcommand {\vect}[1]{\underline{#1}} % tenseur d'ordre 1 de ...
\newcommand {\ten}[1]{\underline{\underline{#1}}} % tenseur d'ordre 2 de ...
\newcommand {\tran} [1] { \stackrel{t}{\rule{0.07cm}{0.cm}
\rule{0.cm}{0.15cm} }\hspace{-0.75ex}#1} % transpose de ...
\newcommand {\tr}[1]{\mbox{Tr}\left( #1\right)} % Trace 

\newcommand {\psc} {.} % produit scalaire
\newcommand {\pdc} {:} % produit doublement contracte

\newcommand {\dete} {\mbox{det}} % determinant
\newcommand {\nab}[1]{\vect{\nabla #1}} % gradient lagrangien
\newcommand {\grad}[1]{\vect{\mbox{grad} #1}} % gradient eulerien
\newcommand {\Div} { \mbox{Div} } % divergence lagrangienne
\newcommand {\dive} { \mbox{div} } % divergence eulerienne

\newcommand {\dsaut} {  [\![ } % saut a gauche 
\newcommand {\fsaut} {  ]\!]} % saut a droite 
\newcommand {\saut} [1] {\dsaut #1 \fsaut} % saut a travers
\newcommand {\rest} [2] {\left. #1 \right|_{#2}} % restriction sur ??
\newcommand {\moy} [1] {\left\langle #1 \right\rangle} % moyenne

%
% DERIVEES
%
\newcommand {\der}[2]{\frac{\partial #1}{\partial #2}} % darivae partielle simple
\newcommand {\deri}[3]{\frac{\partial^{#3} #1}{\partial #2^{#3}}} % iame darivae partielle 
\newcommand {\derd}[3]{\frac{\partial^2 #1}{\partial #2\partial #3}} % darivae partielle croisae d'ordre 2

%
%FONCTIONS PARTICULIERES
%
%\newcommand {\uni} {\ten{\rm 1\hspace{-0.59ex}I}} % Tenseur unita
\newcommand {\uni} {\ten{\mathbbm 1}} % Tenseur unita avec \usepackage{bbm}
\newcommand {\demi} {\frac{1}{2}} % demi (sic!)
\newcommand {\R} {\mathbb R} % espace des raels
\newcommand {\Ci} [1] {\mathscr{C}^{#1}}  % Espace de continuita C_i
%\newcommand {\ca} {\mathscr{CA}}  % Cinamatiquement admissibles
%\newcommand {\sa} {\mathscr{SA}}  % Statiquement admissibles
\newcommand {\ca} {\mbox{CA}}  % Cinamatiquement admissibles
\newcommand {\sa} {\mbox{SA}}  % Statiquement admissibles

\newcommand {\J} {\mbox{J}} % Fonctionnelle

%
% TOUTES COORDONNEES
%
\newcommand {\Ve}  {\vect{e}} % vecteur de base eulerien 
\newcommand {\VE}  {\vect{E}} % vecteur de base lagrangien
\newcommand {\nn} {\vect{n}} % normale a \dDom
\newcommand {\nt} {\vect{t}} % tangente a \dDom
\newcommand {\nb} {\vect{b}} % binormale a \dDom

%
% DOMAINES, FRONTIERES
%

\newcommand {\Dom} {\Omega} % Domaine
\newcommand {\dDom} {\partial\Omega } % bord du domaine
\newcommand {\Do}[1] {\Omega_{#1}} % Domaine  i
\newcommand {\dDo}[1] {\partial\Omega_{#1} } % bord du domaine i
\newcommand {\dDof}[1] {S_{T_#1} } % bord a condition de Neumann en i
\newcommand {\dDou}[1] {S_{\xi_#1} } % bord a condition de Dirichlet en i

%
%QUANTITES MECANIQUES
%
\newcommand {\pl} {\vect{X}} % vecteur position lagrangien
\newcommand {\pe} {\vect{x}} % vecteur position eulerien
\newcommand {\T} {\vect{\Phi}} % transformation du systeme
\newcommand {\F}{\ten F} % gradient de la transformation
\newcommand {\dep} {\vect{\xi}} % champ de deplacements


\newcommand {\C}{\ten C} % dilatations de Cauchy-Green
\newcommand {\e}{\ten \varepsilon} % deformations de Green-Lagrange
\newcommand {\sig}{\ten \sigma} % contraintes de Cauchy
\newcommand {\pio}{\ten \pi} % contraintes de Piola-Kirchhoff
\newcommand {\B}{\ten B} % contraintes de Boussinesq

\newcommand {\K}{\mathbb K} % Matrice de rigidite
\newcommand {\FG}{\mathbb F} % Force generalisee
\newcommand {\CT}{\mathbb C} % Matrice de capacite
\newcommand {\tg}{\mathbb T} % Temperature generalisee

\newcommand {\ener} {w} % energie libre volumique
\newcommand {\Ep} {\mathscr{E}^p} % energie Potentielle totale
\newcommand {\El} {\mathscr{E}} % energie libre totale
\newcommand {\La} {\mathscr{L}} % Lagrangien du systeme
\newcommand {\D} {\mathscr{D}} % Disipation du systeme
\newcommand {\Pe} {\mathscr{P}_e} % Puissance exterieure
\newcommand {\Pui} {\mathscr{P}_i} % Puissance interieure
\newcommand {\W} {\mbox{W}} % energie elastique

%
% VITESSES, TAUX DE QUANTITES MECANIQUES
%
\newcommand {\vit}[1] {\dot{#1}} % vitesse de ...
\newcommand {\acc}[1] {\ddot{#1}} % acceleration de ...
\newcommand {\vdep} {\vit{\dep}} % champ de vitesse
\newcommand {\adep} {\vect{\gamma}} % champ d'acceleration
\newcommand {\vd} {\dep^*} % champ de vitesse CA

%
% ELEMENTS DIFFERENTIELS
%
\newcommand {\dv} {\,d\Omega} % volumique
\newcommand {\ds} {\,ds} % surfacique eulerienne
\newcommand {\dS} {\,dS} % surfacique lagrangienne
\newcommand {\da} {\,da} % surfacique
\newcommand {\dl} {\,d\ell} % lineique
\newcommand {\dz} {\,d\z} % en z
\newcommand {\dt} {\,dt} % en temps

%
% INTEGRALES
%
\newcommand {\intv} { \int_{\Dom} } % sur le domaine
\newcommand {\ints} { \int_{\dDom} } % sur le bord du domaine
\newcommand {\intvi}[1] { \int_{\Do{#1}}} % sur le Domaine  i
\newcommand {\intsi}[1] { \int_{\dDo{#1}}} %  sur le bord du Domaine  i
\newcommand {\intsu} { \int_{\dDou i} } % sur le dommane en Dirichlet dans la direction i
\newcommand {\intst} { \int_{\dDof i} } % sur le dommane en Neumann dans la direction i

%
% Palette de couleur
%
\include{New_BD_ColorScheme}

%%%%%%%%%%%%%%%%%%%%%%%%%%%%%%%%%%%%%%%%%%%%%%%%%%%%%%%%%%%%%%%%%%%%%%%%%%%%%%%%%%%%%%%%%%%%%%%%%%

%\newcommand {} {}

%